\section{The Defeat and the Future of France}

The following essay by \textbf{Julius Evola} was published as ``La disfatta e il futuro dell Francia secondo l'Action Française'' in \textit{La Vita Italiana}, in April 1942. Due to its length, it will be posted in two installments.

\begin{quotationx}
This essay may be of historical interest to some, since, given that history is written by the victors, it offers a rare glimpse into the thinking of the losing sides in World War II. Evola is reviewing a book by Maurras that deals with the quick and humiliating defeat of France by the Germans.

There are recurring themes that repeat today. First of all, Maurras distinguishes between the real and the legal countries\footnote{\url{https://en.wikipedia.org/wiki/Maurrassisme\#A\_political\_model\_based\_on\_the\_.22Real\_Country.22\_.28pays\_r.C3.A9el.29}}; the first is organic, consisting of those elements who embody the soul of France and the latter a strictly legal entity comprised of all those inhabiting a geographic area. Maurras derives many of his concepts from Medieval social theory. If contemporary Catholic are embarrassed by this, it is only because they are preaching a different religion.

Maurras points to partisan interests and party politics as a cause of weakness. We will simply point out, as does Maurras\footnote{\url{https://gornahoor.net/?p=5247}}, that Right without the force to back it up is worthless. There are advantages to winning and consequences to losing. In the second part of this essay, Evola will develop his own critique of Maurras.
\end{quotationx}

\paragraph{First part}
In a recently published book that quickly reached its 29\textsuperscript{th} edition, Charles Maurras sought to penetrate the meaning and the causes of the tragic destiny that weighed on his country if only to indicate the ways along which France will be able in the future to bring itself back and restore itself, drawing solutions from the hard trial. \emph{La seule France, Chronique des jours d'épreuve}, France alone, Chronicle of the days of trial, 1941 We believe it is not lacking interest, for our readers, to reproduce the principle theses of this book, given the personality of its author.

First of all, about the causes of the disaster. There are internal and external causes. Regarding the internal causes, Maurras flatly opposes the thesis of those who speak of a morally degraded France, physically degenerated, and decayed military virtues. The country of Joan of Arc, who already overcame the trial of 1914, should have been able to demonstrate, not less than any other nation, Germany included, the possibilities of a national exaltation of a heroic impulse. About this last point, in a special chapter, Maurras remembers, on the contrary, a series of cases of warrior and heroic virtues of the French army, recognized, in this war, by the enemy himself, but demonstrated sporadically and in disparate circumstances, through the lack of all the conditions of a not unequal battle.

After making the distinction between the \emph{real country} and the \emph{legal country}, Maurras instead believes that he indicates the principle internal cause of the French defeat in the split and, then, the division and denationalization between these two countries. The natural resources of France will be divided through the split nature of its own government. While in Germany, a political factor that united and multiplied every force had the upper hand, in France an similar factor served only to break it up: power was in the hand of politicians dedicated to parliamentary discussions and chaotic successions: fifty ministers in twenty years and one hundred ten in seventy years. Moreover, this government was something superimposed on the real country, having in view mostly the interests of the part and the party---apart from the influences of international powers---than those of the nation. So while in Germany there was an intensive culture of patriotism, in France it was always more neglected and rhetorical. While in Germany there was an improvement and a concentration of all the forces of the State, in France there was deliquescence and dislocation.

Maurras was opposed to the idea that authoritarianism would be something German:

\begin{quotex}
In the country of Louis XIV and Richelieu the concentration of the State would have been equally conceived and pursued according to a model of the most straightforward and most original French method.
\end{quotex}

The principle internal cause of the French disaster would therefore be in that. But, to tell the truth, Maurras forgot that the French recovery of 1914 happened even without such conditions, having instead democratic and parliamentary France as its antecedents. Therefore other causes are also looked for: that Maurras forebodes when he speaks of the ``constructed'' character of the French war and indirect powers that determined it, by means of a type of hypnosis exercised on public opinion, after all those national elements that would have been able to oppose the manoeuvre were paralyzed. According to Maurras, it would have been up to the French nationalists to prevent the war in 1935, during the period of sanctions, then in 1936-1937 on the occasion of Spain and in September 1938, on the occasion of the Czechoslovakian crisis. But in August 1939 all the measures were taken to prevent public opinion from being influenced in a direction contrary to what was wanted: seven or eight days before the declaration of war a decree was issued that ``blocked'' the press and authorized the immediate suppression of all these newspapers that intended to oppose the warmongering tendencies. ``What is visible in this war,'' Maurras writes:

\begin{quotex}
... is the constant progress of a plan, the development of a patiently hatched plot, to attract a herd of idiots in the trap constructed by the wicked... the declaration of war of September 3 was obtained on the basis of everything there is in the \emph{demos} that is soft, vain, empty, absent, and nonexistent; it had to cross---and it crossed---the empty spaces of hesitation, imbecility, incoherence, and a more than animal stupidity: mineral---that of the billiard ball and whatever other heavy body provided with any mobility.
\end{quotex}


If the origin of such a declaration of war

\begin{quotex}
was covered by vapours so sooty, it is fated that it was so rapidly dissipated.
\end{quotex}


Then, the lack of reaction of the national French soul after the first moment and, especially, after the first period of the listless and only nominal war on the Maginot line. If we could oppose to that, that the awakening would nevertheless have had to be manifested once that, France being at this point at war, its national territory was directly threatened, Maurras indicates, here, the other technical and external causes relative to the French military unpreparedness and English treachery.

For example, he points out a typical case of deceit to the country. In the Chamber, with respect to aerial armaments, an engineer made a display showing that against 4000 German aeroplanes, the French could oppose only 800 and almost all of them fighters. The Head of the Government and the Minister of Air rushed to deny this assertion, demonstrating that the French air potential was instead about 2300 aircraft---making use however of inexact data, and referring simply to ``naked'' aeroplanes, far from being functional. Moreover, it counted on the help of the ``English aviation power''. This is one of the many cases of the blindness and irresponsibility that, from the technical side, had a great weight in the French catastrophe, so that Maurras speaks of a true betrayal of the leaders vis-a-vis the French nation, except the remaining military honor.

Maurras held the idea that the only path that an intelligent and national government in France should have taken, especially after the events of 1938, was that of Rome. He accused France of not having followed the real national interests, but of being left to take part in a mere ideological war. It should have been clear that France was not able to do anything at all for Poland. Also suddenly initiating a serious war on the Maginot line, that would have been, says Maurras, to try to break through a wall with one's head to help someone who was being assassinated outside of it. Other accusations were made in regard to the advance in Belgium, something that again would have obeyed suggestions of the English, concerned to maintain the coasts of the English Channel, more than the true strategic reasons for the French defense.

On all these points we cannot therefore be fully in agreement with Maurras. The way toward Rome, for France, would have been possible only in the supposition of a preliminary upheavel in the national sense, at least to not reduce the whole to a simple compromise with its precise mental reservations. But in spite the efforts of Action Française and similar organizations, in France there was not ever a trace of such an upheaval, the predominant mentality being in fact far from the ideas of Fascist renewal. So as things stand, one cannot dispute that those who wanted to give to the war an essentially ideological character are not right---from their point of view, naturally. It is useless to hide it: sooner or later, with the front line of the democracies, plutocracies, and Bolshevism, it should have come to daggers. A new Monaco in 1939 would have only meant a new respite. Given the direction of the dynamism of the ``fascist'' nations, we believe it is difficult that it would have been able to reach a true and durable pattern of balance without a violent solution and to give Europe a new shape.

The intervention of France, therefore, in my opinion, returns perfectly to the logic of those who were leading it, at least as to the legal country, if not the real country. From such a point of view, a Frenchman should instead have accused the indefensible and incomprehensible blunders that the democratic front committed, beginning with Versailles, in the sense of not taking all the measures to thereby prevent the enemies from being revived. Instead of taking it, for ideological reasons and essentially for masonic and anti-traditional hatred against the Hapsburg Empire, it should have sought to dismember Germany and, at the time of the occupation of the Ruhr, to act without hesitation. Step by step, the demo-masonic front instead allowed the enemy to gather strength. Therefore, the only reasonable accusation to make, instead, would be that which, after arriving at 1939 and having ``swallowed'' the annexation of Austria, the division of Czechoslovakia, the weakening of the Little Entente\footnote{\url{https://en.wikipedia.org/wiki/Little\_Entente}}, and the definitive strengthening of the Axis Powers, it was by that time too late to do anything with the assurance of success: it was rather, especially for France, heading for suicide, because of the instinctive reaction awakened by an anxiety complex.

\paragraph{Second part}

This is Part II of the essay by \textbf{Julius Evola} that was published as ``La disfatta e il futuro dell Francia secondo l'Action Française'' in La Vita Italiana, in April 1942.

\begin{quotationx}
This piece is of historical interest since the issues in 1941 are still in play in 2013. We also get an ``inside'' look at certain thought processes when both Evola and Maurras expected victory for their side. Philippe Pétain and Charles Maurras had been well respected figures in France before the war but then ended their days in disgrace.

The scission in France referred to was between two generals: Charles de Gaulle and Petain, the former representing the Republic and cosmopolitanism and the latter, identitarianism. I don't know the specific features of the laws against the Jews; unfortunately, neither Evola nor Maurras admits that France was under the thumb of Germany, so that France's internal ``Jewish problem'' was more far-reaching as many were sent off to concentration camps. This policy resulted in the discrediting of identitarian movements to this day as they stoke up that very real fear.

Defeat has consequences, so the victory of the Allies has established republican and cosmopolitan values as the default position of the West. Jewish and Masonic values are now considered part of the mainstream of Western thought and beyond debate. The third element, that of foreign immigration, is still contestable, although the opposition is weak and unconvincing.

The words must have stuck in his throat, since Evola could not quite name the ``spiritual and traditional authority''. On the other hand, I doubt anyone in the Vatican II church would be happy to know that the Catholic Church was at the spiritual center of Vichy France. This shows how much the spirit of the church has changed in the intervening 70 years, since it was then considered compatible with nationalism. Although it is now more cosmopolitan and its sophisticated social teachings are now little more than ``Marxism light'', this has not led to more acclaim by the modern world as the architects of Vatican II had hoped, but instead has led to its marginalization.
\end{quotationx}


However, that does not prevent what Maurras documents and says about English politics, which had not the least regard for French interests, from being completely fair. Maurras takes a truncated position against De Gaulle. He, on the contrary, dedicates a chapter to demonstrate that England, basically, continued to be even in more recent times, in spite of appearances, the enemy of France, jealous of its colonial empire: it has always done everything to impede them from taking the rank that was due to them. Maurras fears that England seeks to exploit the French disaster to realize definitively their hidden ambitions, indeed making use from De Gaulle for this end, and stoking the internal French schism by every means.

For this reason, Maurras exhorts his compatriots to convince themselves of the absolute necessity of an internal concentration on the government of Marshall Petain as conditions for any recovery of France, after such a hard and tragic experience. So he devotes the second part of his book to highlighting and developing everything that the new Government tried to do. Maurras, although a realist, is pro Petain, because he eventually moved to establish a personal, national, authoritarian government, saved from parliamentary bedlam. France must, in that way, rediscover the true idea of the State: the State connected to the real country, freed from international influences, strong and disciplined, but against every leveling centralization.

Having such a State in mind, Petain had already begun to act to return France to the French: and here Maurras illustrated the new laws intended to build a dam against the influx of foreigners, metics, and immigrants that previously had found in France a type of promised land. Laws against Jews and Masons followed. Regarding the Jews, Maurras made the French point of view clear: they do not originate from an anti-semitic hatred, but from the fact that a Jew can only be considered to be an alien, who as such cannot be equated to the citizens of a country and moreover must accept the conditions that are imposed on whoever desires to be a guest. Maurras however correctly observes that French legislation on the Jews conveys mainly a juridical point of view, inspired by effects more that by causes: Maurras' central point is to understand that a style and a hereditary mentality, a tradition and indelible custom will always make the Jewish substance something inassimilable and corrosive. Maurras recalls that France seized a large number of significant Jewish refugees, starting with the five Rothschilds. It advanced a sensible proposition: to appropriate such fortunes for the restructuring of rural debt, i.e., of those contracted or to contract for the acquisition and cultivation of land. That would mean to strengthen the class which more must be counted on for the reconstruction of a nation. As to the anti-masonic battle, Maurras brings to light the danger inherent in proceeding, in this regard, with a direct and practical action, extremely difficult where it is about a secret association. They can have effective results when they have complete possession of the State. As Mussolini did, only then will it be able to little by little paralyze the Masonic influence. The first step, i.e., the formal prohibition of Masonry, was already accomplished by France. What followed is:

\begin{quotex}
to promote a moral publicity without rest and carry everywhere the tip of its fire.
\end{quotex}


Maurras then alludes to social reforms: the question naturally arises the reconciliation of classes and social justice. But, in this regard, Maurras also makes a rather appropriate point:

\begin{quotex}
If the protection of the poor is the natural duty of the State---when is not exercised spontaneously by some powers worthy of their strength---it must however also defend isolated, but just, powers against certain coalitions of ``the poor'' who gain an enormous voice, capable of crushing everything under the darkest and most brutal tyranny.
\end{quotex}


And France, which first had the Revolution (revolution, says Maurras correctly---is still so stupid to be called French, being instead the most cosmopolitan of all of them), has from historical experience everything that it needs in order to understand, better than any other nation, the timeliness of a similar warning.

After some other specific points, Maurras addresses the reform of education in France, indicating what was already done and what still is to be done. The fundamental need is that the State in France is made truly father of its schools by proceeding systematically in the formation of a new mentality in the next generations. In our opinion, this is really the crucial point, on which almost everything else depends. The French mentality still needs to be detoxified from the deleterious effects of ideology, which in that country has existed for a while and had a characteristic diffusion, in width as in depth. And that, in good measure, must happen through the schools.

Maurras indicates here correctly a point that, in the current state of reform, was not sufficiently considered; i.e., that which is not limited to a more or less autonomous moral teaching, but at whose center the principle of a spiritual and traditional authority is returned, the only one that can make well understood all the errors and deviations in which France has incurred: the return to the classical spirit, the battle against romanticism and humanitarianism, the decisive antithesis to the Revolution and all the phenomena of degeneration, decadence, and anarchy that were its consequences.

In one point of his book Maurras develops some rather important considerations in connection to the idea that the French disaster of 1940 seems to be the ultimate effect of a chain of causes that go back far in time and of that only some are of general rule.

\begin{quotex}
The mass that constitutes an avalanche is formed of sand and stone coming from the most various rocks, something that however does not prevent the observer from understanding the small movement or the smallest human push that detached the first lump of snow---thus, before the evidence of any accusing traces (an impression of human finger), before the still visible trajectory of any suspicious footprint (explicitly oriented, clearly directed), the watchful eye discovers in the right point the impulse and the form of a criminal act---that waits only to be identified with the name of its author.
\end{quotex}


This is an observation which demonstrates, in Maurras, the right orientation to truly deepen the secret history of his country and therefore even to identify the premises of a true reconstruction. Here we are in the area of what we call the ``science of subversion'': in regard to which perhaps no country in the world offers such precious research material, as does France.

If we have to indicate any defect in Maurras' book, it is that it does not touch at all the problem of the new European order and the corresponding supernational arrangement. We can well understand, in Maurras, the need connected to the formula: \emph{France alone first}---he himself indicates as much its provisional nature, when he says that, if there is going to be a European ``concert'', it is first necessary to make France able to play its part in it. But this image makes clear exactly the gap noted: a single part can be played correctly only by taking the movements from the central motif of an entire symphony. We therefore propose to French nationalism the further task of also taking this point into account and proceeding to an action of mental formation in the sense, as soon as the more immediate needs of reconstruction and internal recovery are realized.

\flrightit{Posted on 2012-11-19 by Aeneas}

\begin{center}* * *\end{center}

\begin{footnotesize}
\begin{sffamily}
\texttt{Logres on 2012-11-20 at 12:47 said:}

De Gaulle gained his authority from being in charge of one of the few military counter-successes against the Germans, which actually happened at a few points -- French arms were not inferior to Germans, whose military equipment was not at all overwhelmingly superior. Does Maurras have thoughts about De Gaulle? France is still a strange country -- on the one hand, they banned the burka, but on the other, seem willing to tolerate the multicultural hordes seeping into their body politic today. One can't help but think that the Allies were punished for their open support of the Red cause in Spain during that terrible civil war, for spreading bloodshed amongst the Catholics, who were really fighting for their lives (although Franco wasn't exactly a solar king).

\hfill

\texttt{Jason-Adam on 2012-11-20 at 17:53 said:}

This is indeed a very interesting piece and I thank Cologero for the translation.

In reading it my symapthies tend toward Maurras' thesis and not Evola's criticisms... this line struck me ``Given the direction of the dynamism of the ``fascist'' nations, we believe it is difficult that it would have been able to reach a true and durable pattern of balance without a violent solution and to give Europe a new shape.'' Is Evola advocating agressive warfare ? Not only is intra-Christian warfare morally wrong, it's also stupid when Russia was looming to swoop in and enslave all to pick fights over modern borders. If he really wanted a heroic sacred war, a new crusade, why not call on all Europeans and men of Tradition to unite as one and take up arms to destroy Communism ? The success of Spain \& Portugal showed that neutrality was an option for France. An aquaintance of mine who is very much in the know told me once that if Italy had stayed neutral, she would have emerged the strongest nation in Europe after the war and fascism would have been extremely attractive then to all the ruined nations... ... .the way I see it, allowing the Protestant nations to blow themselves up while the Catholic nations rested in peace would have made a far more frutiful world from a Traditional standpoint.

Franco was a hero, the living example of a dictator in Donoso Cortes's thoughts -- a man who acts with power to hold back the tide.

What Catholics are criticising Gornahoor ? I haven't seen anything here that contradicts the magisterium of the Church. If you mean Modernists you must remember that they are not Catholics, they are heretics and condemned by Catholic dogma. Even if he wears a mitre if a man teach error he is a heretic and a true Catholic must stand up to defend the Faith.

\hfill

\texttt{Logres on 2012-11-20 at 18:46 said:}

Several Catholics have argued \& then left the site over what can only be described, the best I could tell, as a refusal of Cologero to speak certain phrases, or refrain from making other suggestions. This despite all efforts being made to accommodate. But that's my perspective, \& I didn't follow it closely.

\hfill

\texttt{Jason-Adam on 2012-11-20 at 19:05 said:}

I can not speak for others, but I am a Catholic and find nothing in Gornahoor that opposes the Faith.

\hfill

\texttt{Cologero on 2012-11-20 at 23:42 said:}

De Gaulle was earlier associated with Action Francaise. Obviously, the war separated them, Maurras becoming a collaborator and de Gaulle the leader of the Free French.

There is an important point made in this piece, and its consequences will be clear in the second part. Was the war to serve French national interests or was it an ideological war. Petain, and Maurras supported him, saw it the first way and implemented policies whose purpose was to promote and protect French interests. However, in retrospect, we see WW II as an ideological battle, viz., as a war of the free democracies (good) against totalitarian Fascism (evil). This view still dominates the self-understanding of the West to this day. Thus, we will see that Petain's policies were reprehensible by today's standards. Mircea Eliade reminds us of the eternal return and we will how this ideological war is still being played out in Europe and North America. To understand this essay will be to understand what is happening today.

\hfill

\texttt{Cassiodorus on 2013-02-03 at 14:30 said:}

Off topic... Has anyone seen the recent debate between William Lane Craig and the atheist philosopher Alex Rosenberg at Purdue university?

It seems to me that Craig soundly defeated yet another representative of our free-thinking sophisticated intellectual elite.

\hfill

\texttt{Cologero on 2013-02-04 at 12:39 said:}

Will all due respect, Cassiodorus, I'd prefer to stay on topic, although I see there is little interest in this topic.

As for that ``debate'', we can't even accept its premises. There is a distinction between philosophy (or what counts for it now) and metaphysics (at least what it used to mean). The former like ``debate'', the later prefers knowledge. So it is not a question of proving a ``reasonable belief in God, but rather in ``knowing'' God, as the catechism says.

Keep in mind that there can be no ``rules'' for a debate such as those proposed, other than to be logical and keep your facts straight. A termination point must be agreed to otherwise there is endless chatter. Decision is necessary. Has any participant in similar debates suddenly ``resigned'' by being convinced by his opponent? No, these debate are more like chariot races with their fans taking sides before the first word is even spoken.

Another point is that the participants make a career out of such debates; that makes them more like sophists than true philosophers. As for Mr. Craig, we don't accept the existence of the being he calls ``God'' anyway (at least as far as I understand him), so I don't know how he can ``win''. Consult \textit{this}: \url{https://edwardfeser.blogspot.com/2009/11/william-lane-craig-on-divine-simplicity.html}, for example.

\hfill

\texttt{Jason-Adam on 2013-02-04 at 13:25 said:}

I find it interesting that Evola finds himself in accord ultimately with Maurras and he does see nationalism in a positive light as a necessary first step in the battle against subversion. Quite a contrast from de Benoist's critique of nationalism which while correct vis a vis the concept of the empire winds up supporting the communist destruction of all identity.

\hfill

\texttt{Cologero on 2013-02-04 at 17:50 said:}

True, Jason-Adam, Evola's only objection is that Maurras did not address the question of a super-national arrangement, although it is not clear to me what Evola himself may have had in mind. Perhaps Evola was not familiar with Maurras' 1922 essay on Latinity which did address that question.

I really am not familiar with Benoist's critique. What I don't see among the neo-pagans is how they address the issue of the spiritual and traditional authority which must be at the heart of any future return to Tradition; this would order the socio-economic systems. Does Benoist have a complete socio-economic position of his own?

\hfill

\texttt{Jason-Adam on 2013-02-04 at 20:52 said:}

De Benoist has said that he is not a Traditionalist in the sense of believing in a Philosophia Perennis so it iz puzzling that anyone could see him as a successor to Guenon or Evola.

He supports direct democracy and has stated that democracy is part of the Greco-Roman pagan tradition which Christianity destroyed so puts him right in line with the anti-mediaeval critique of the modern world and against the aristocratic conservatism of a de Maistre or Cortes which Evola supported.

I think the main similarity between Evola and de Benoist is simply their Germanophilia, other than that they are widely divergent thinkers.

\hfill

\texttt{Logres on 2013-02-04 at 21:12 said:}

It's interesting that the neo-pagans and the proto-Fascists, despite having abundant intellectual energies and keen insight into certain aspects of Modernity, as well as having a will to resist, have not been able to coordinate among themselves an effective consensus, except at the highest intellectual levels (Evola, Maurras, etc., who all read each one another), and that decades ago. It's almost as if the impetuous rush towards Chaos in the energies of Modernism has infected the adherents of these insightful and courageous intellectual giants. I really need to find more translations of Maurras in English. Especially the essay on Latinity (which is probably not available in ghetto-English). It might be worthwhile to define the ``Classicism'' of France (and also Europe) as something over against the supposed ``Classicism'' (or Republicanism/Grecianism) that dominated the events of 1789. Evola seems to allude to that in his ``material for study'' remark concerning those events.

\hfill

\texttt{Cassiodorus on 2013-02-05 at 16:27 said:}

Cologero,

I basically agree with you. I am certainly no proponent of theistic personalism; It's definitely the classical theism of the Thomist for me. As far as the ``debates'' go, I have mixed feelings. Of course, you're quite right about the ``chariot races''. Nevertheless, I think Craig's efforts can bear some good fruit. To hear an analytical ``philosopher'' effectively challenge the ``default'' modern worldview of materialism and relativism in large public venues could be persuasive for our largely agnostic culture.

\hfill

\texttt{Cologero on 2013-02-06 at 07:33 said:}

Maybe you know something I don't, but I don't see these debates persuading anyone. There may be some poaching at the margins in both directions, but no large scale movements. To change the default modern worldview would require a change in body biochemistry which is more than an argument is capable of. Now all Craig tries to do is to show that belief in his god is ``reasonable''. However, I don't think his belief is reasonable at all ... it is too close to a ``Mr God'' or even to the Mormon conception.

What may be more useful is a series of debates between Feser and Craig ... effort is better spent at this time, I think, to reduce the number of heretics rather than to convert the heathens.

\hfill

\texttt{Cassiodorus on 2013-02-06 at 11:37 said:}

A series of debates between Craig and Feser- Yes! That would be be excellent. In fact, I think some folks have even suggested that on Feser's blog. Cologero, I have taken note that the Aristotelian-Thomistists practically short circuit when the Divine Principle of ``beyond-being'' is brought into the conversation. I've been reading David Bentley Hart lately in an effort to understand the exact nature of the ``dispute'' between the Christian West and East.

\hfill

\texttt{Cassiodorus on 2013-02-06 at 11:40 said:}

Did I say, Thomistists? Wow.

\hfill

\texttt{Cologero on 2013-02-06 at 23:06 said:}

Cassiodorus, perhaps you can comment on this: \textit{Orthodoxy and Aquinas}\footnote{\url{https://easternchristianbooks.blogspot.com/2013/02/marcus-plested-on-orthodoxy-and-aquinas.html}}

\hfill

\texttt{Cassiodorus on 2013-02-07 at 14:50 said:}

I'll check it out.

Many thanks!!!

\hfill

\texttt{Jason-Adam on 2013-02-08 at 15:08 said:}

The book on Orthodoxy and Aquinas is worthy of reading, it is quite the case that in the Middle Ages the eastern church was not so hostile towards Aristotle and scholasticism as one would expect, even the arch-anti-west Patriarch Gennadius was a strict Scholastic and hated the neo-platonism of Plethon.

This proves to me more and more that we can not understand developments in philosophy without understanding geopolitics. The Orthodox Church in Russia had to become anti-west in theology to justify Russia's sense of non-Europeanness and unique culture that is opposition to the west.

Which is why I question that, even if somehow Europe returned to Tradition, there would still be enmity between the 1st and 3rd Romes until one prevails.


\hfill

\end{sffamily}
\end{footnotesize}

